\begin{abstract}
Smart warehouses need wireless networks that can support sensors, controllers,
video traffic, and mobile robots while also reacting to changes in the
environment. This paper presents a situation-aware Integrated Sensing and
Communication (ISAC) co-simulation framework that connects Sionna RT ray
tracing with ns-3/5G~LENA networking. The framework uses ray-traced multipath
to detect reflected targets, filters clutter, tracks moving objects, and then
uses the confirmed target positions to form sensing-informed beams.

The framework is evaluated in two stages. First, controlled free-space and
warehouse scenes are used to test the sensing and beamforming pipeline under
clear and cluttered conditions. Second, an end-to-end warehouse Digital Twin is
used to study the same mechanism with mobile robots, package and rack sensors,
cameras, video clients, and mixed application traffic. The controlled results
show that ISAC reduces path loss in both scenes, that a wider beam gives the
best communication behavior, and that tracking is needed to suppress false
detections. Deeper propagation gives a small detection gain, but it can also
increase clutter in the warehouse. In the full warehouse use case, the
$4{\times}4$ gNB array gives the best balance: the $2{\times}2$ array is not
stable enough, while the $8{\times}8$ array increases power consumption without
a proportional performance gain. These results show that practical ISAC
evaluation must consider sensing quality, communication performance, and energy
cost together.
\end{abstract}

\begin{IEEEkeywords}
Integrated sensing and communication (ISAC), 6G, situation awareness, ray
tracing, Sionna RT, ns-3, 5G LENA, adaptive beamforming, multi-target tracking,
Digital Twin, smart warehouse.
\end{IEEEkeywords}
